\chapter{Ricardian Trade in Full}\label{ch:b3:ricardian-trade}

In 1859, after more than two centuries in which almost no foreign trade was
allowed, Japan opened its ports. Within a few years it was exporting silk,
silkworm eggs and tea, and importing cotton yarn, cotton cloth, woollens, rice
and sugar. % ledger: b3rt:bb04.t1
For an economist it was close to a laboratory: a country whose prices had been
set by its own costs alone was suddenly exposed to world prices. \omterm{def:b3:ricardian-trade:comparative-advantage}{Comparative
advantage} predicts which way the goods should flow. The school volumes stated
the prediction for two goods and two countries. This chapter states it for any
number of goods, finds the wage that settles who makes what, adds the cost of
shipping, asks who gains when a trading partner becomes more productive, and
looks at the tests.

\begin{openingchart}{Japan's trade in each year from 1868 to 1875, valued at
the prices of 1851--53, before the ports opened (millions of ryō; the ryō was
worth about one dollar in 1873). Bars: imports (negative) and exports (positive),
each the sum of the components printed in the source; dots: the printed total.
The law of comparative advantage predicts a negative total.}{Bernhofen and Brown
(2004), Table 2; sums of components ours}
\begin{tikzpicture}
\begin{axis}[omtime, width=0.8\linewidth, height=0.42\linewidth, xmin=1867.4, xmax=1875.6,
  ymin=-14, ymax=7, xtick={1868,...,1875}, xticklabel style={/pgf/number format/1000 sep={}},
  ylabel={Millions of ryō at 1851--53 prices}, legend columns=3,
  legend style={at={(0.5,1.03)}, anchor=south}]
  \addplot[ybar, bar shift=0pt, bar width=12pt, area legend, fill=omDef!25, draw=omDef!60]
    table[x=year, y=imports] {figdata/bachelor-3/07-ricardian-trade/out/japan.csv};
  \addplot[ybar, bar shift=0pt, bar width=12pt, area legend, fill=omThm!45, draw=omThm]
    table[x=year, y=exports] {figdata/bachelor-3/07-ricardian-trade/out/japan.csv};
  \addplot[only marks, mark=*, mark size=2.2pt, black]
    table[x=year, y=total] {figdata/bachelor-3/07-ricardian-trade/out/japan.csv};
  \draw[guide] (axis cs:1867.4,0) -- (axis cs:1875.6,0);
  \legend{imports, exports, total}
\end{axis}
\end{tikzpicture}
\end{openingchart}

Each year Japan's imports, valued at what they would have cost in Japan
before trade, were worth more than its exports valued the same way. The
total is negative in all eight years, from $-0.18$ million ryō in 1868 to $-8.06$
million in 1875. % figdata: japan
(The printed totals differ from the sums of the printed components by at most
0.02, an effect of rounding.) In words: Japan sold what had been cheap at home
and bought what had been dear. \Cref{prop:b3:ricardian-trade:law} below shows
why the sign is the test.

\begin{recall}
The grades 10--12 volume taught Ricardo's model with two goods and two
countries: opportunity cost, the production possibility frontier,
specialisation and the \omterm{def:b3:ricardian-trade:terms-of-trade}{terms of trade} between the two \omterm{def:b3:ricardian-trade:autarky}{autarky prices}. The Year 2
volume built general equilibrium with the Edgeworth box. Here the same logic is
extended to many goods and closed with the labour market.
\end{recall}

\begin{notation}[Symbols of this chapter]\label{not:b3:ricardian-trade:symbols}
Home and Foreign (starred) have labour forces $L$ and $L^*$ and wages $w$ and
$w^*$; $\omega = w/w^*$ is the relative wage. Good $i$ (or $z$, on a continuum)
needs $a_i$ hours of labour at Home and $a_i^*$ abroad; $A_i = a_i^*/a_i$. $b(z)$ is
the share of spending on good $z$ and $\vartheta(z)$ the share on goods up to $z$.
$\kappa \ge 1$ is an \omterm{def:b3:ricardian-trade:iceberg}{iceberg} factor (not a regression coefficient). $T$ is a vector
of net exports and $p^A$ of \omterm{def:b3:ricardian-trade:autarky}{autarky prices}; $\lambda$ is a \omterm{def:b3:ricardian-trade:acr}{domestic expenditure
share} and $\varepsilon > 0$ a trade elasticity.
\end{notation}

\section{From two goods to many}\label{sec:b3:ricardian-trade:chain}

\begin{definition}[Autarky price]\label{def:b3:ricardian-trade:autarky}
\index{autarky price}
The \emph{autarky price} of a good is its price in a country that does not
trade. With labour the only factor and constant labour requirements, the
autarky price of good $i$ at Home is $w a_i$, and relative autarky prices are
ratios of labour requirements: $p^A_i/p^A_j = a_i/a_j$.
\end{definition}

\begin{definition}[Comparative advantage; chain of comparative advantage]\label{def:b3:ricardian-trade:comparative-advantage}
\index{comparative advantage}\index{chain of comparative advantage}
Home has a \emph{comparative advantage} in good $i$ relative to good $j$ if its
relative \omterm{def:b3:ricardian-trade:autarky}{autarky price} of $i$ is lower than Foreign's: $a_i/a_j < a_i^*/a_j^*$, or
$A_i > A_j$. Ranking all goods so that $A_1 > A_2 > \dots > A_n$ gives the
\emph{chain of comparative advantage}: Home's advantage is strongest in good 1
and weakest in good $n$.
\end{definition}

\begin{proposition}[The relative wage cuts the chain]\label{prop:b3:ricardian-trade:cut}
With perfect competition and free trade, at relative wage $\omega$ Home makes
every good with $A_i > \omega$ and Foreign every good with $A_i < \omega$; a good
with $A_i = \omega$ may be made in both.
\end{proposition}

\begin{proof}
Under perfect competition a good's price equals its lowest unit cost, $\min(w
a_i, w^* a_i^*)$, and only the countries achieving it produce. Home's cost is
lower when $w a_i < w^* a_i^*$, that is when $\omega < a_i^*/a_i = A_i$; Foreign's when
$A_i < \omega$; the two are equal when $A_i = \omega$.
\end{proof}

\begin{omfigure}
\begin{tikzpicture}[econ, x=1.05cm, font=\footnotesize]
  \draw[->] (0,0) -- (11.2,0) node[right] {$A_i$ falls};
  \foreach \x/\g/\a in {0.6/grain/6.0, 1.9/timber/4.0, 3.2/fish/3.0, 4.5/wool/2.4,
     5.8/cloth/1.5, 8.2/tools/1.25, 9.5/shoes/1.0, 10.8/radios/0.6} {
    \fill[omDef] (\x,0) circle (2pt);
    \node[above] at (\x,0.1) {\g};
    \node[below] at (\x,-0.1) {$\a$};
  }
  \draw[omThm, very thick] (7.0,-0.9) -- (7.0,1.0) node[above] {$\omega = 1.4$};
  \node[omDef, align=center] at (3.2,-1.1) {made at Home (Brenland)};
  \node[omDef, align=center] at (9.5,-1.1) {made abroad (Kestria)};
\end{tikzpicture}
\par\smallskip{\small The \omterm{def:b3:ricardian-trade:comparative-advantage}{chain of comparative advantage}, with some of the goods
of this chapter's weekend problem: goods ranked by $A_i$, Home's relative
productivity, and cut by the relative wage. Raise $\omega$ and the cut moves left:
Home makes fewer goods.}
\end{omfigure}

The proposition settles who makes what once $\omega$ is known. It does not say
what $\omega$ is; that needs the demand side and the labour market.

\section{A continuum of goods}\label{sec:b3:ricardian-trade:dfs}

\begin{definition}[Dornbusch--Fischer--Samuelson model]\label{def:b3:ricardian-trade:dfs}
\index{Dornbusch--Fischer--Samuelson model}
The \emph{Dornbusch--Fischer--Samuelson model} has two countries, one factor,
and a continuum of goods $z \in [0, 1]$ ranked so that $A(z) = a^*(z)/a(z)$ is
continuous and strictly decreasing. Consumers everywhere spend a fixed share
$b(z)$ of income on good $z$ (Cobb--Douglas preferences), with $\vartheta(z) =
\int_0^z b(s)\, ds$. Trade is free and balanced, and labour is fully employed.
\end{definition}

If Home makes goods $[0, \bar z]$, world spending on Home's goods is $\vartheta(\bar
z)$ times world income, and it must pay Home's workers: $wL = \vartheta(\bar z)(wL +
w^*L^*)$. Solving for the relative wage,
\begin{equation}\label{eq:b3:rt:B}
  \omega = \frac{\vartheta(\bar z)}{1 - \vartheta(\bar z)}\,\frac{L^*}{L} \equiv B(\bar z),
\end{equation}
which rises with $\bar z$: the more goods Home makes, the more of world spending
falls on its labour, and the higher its wage must be. \Cref{prop:b3:ricardian-trade:cut}
gives the other relation, $\omega = A(\bar z)$.

\begin{proposition}[The DFS equilibrium exists and is unique]\label{prop:b3:ricardian-trade:dfs}
If $b(z) > 0$ for all $z$, there is exactly one $\bar z \in (0, 1)$ with $A(\bar z) =
B(\bar z)$, and the equilibrium relative wage is $\omega = A(\bar z)$.
\end{proposition}

\begin{proof}
$B$ is continuous and strictly increasing, since $\vartheta$ is, with $B(0) = 0$ and
$B(z) \to \infty$ as $z \to 1$. $A$ is continuous, strictly decreasing and positive.
The difference $A - B$ is therefore continuous and strictly decreasing, positive
at $z = 0$ and negative near $z = 1$; by the intermediate value theorem it is
zero at some $\bar z$, and strict monotonicity makes that zero unique.
\end{proof}

\begin{omchart}{Simulated: the Dornbusch--Fischer--Samuelson diagram for 200
goods with equal spending shares and a Foreign labour force 1.5 times Home's.
$A(z)$ is Home's relative productivity, falling along the chain; $B(z)$ is the
relative wage at which trade balances when Home makes goods up to $z$.}{Simulated
by the authors (seed 1977)}
\begin{tikzpicture}
\begin{axis}[omtime, width=0.72\linewidth, height=0.42\linewidth, xmin=0, xmax=1, ymin=0,
  ymax=4, xlabel={Good $z$}, ylabel={Relative wage $\omega$},
  x tick label style={/pgf/number format/.cd, fixed, precision=1}]
  \addplot[omDef, thick, no marks, const plot] table[x=z, y=A] {figdata/bachelor-3/07-ricardian-trade/out/dfs_A.csv};
  \addplot[omThm, thick, no marks, sharp plot] table[x=z, y=B] {figdata/bachelor-3/07-ricardian-trade/out/dfs_B.csv};
  \draw[guide] (axis cs:0.4475,0) -- (axis cs:0.4475,1.2106);
  \draw[guide] (axis cs:0,1.2106) -- (axis cs:0.4475,1.2106);
  \node[font=\scriptsize, anchor=south west, omDef] at (axis cs:0.05,2.9) {$A(z)$};
  \node[font=\scriptsize, anchor=west, omThm] at (axis cs:0.74,3.6) {$B(z)$};
  \node[font=\scriptsize, anchor=north east] at (axis cs:0.44,1.15) {$(\bar z, \omega)$};
\end{axis}
\end{tikzpicture}
\end{omchart}

In the example Home makes the first 89 of the 200 goods and earns 1.21 times
Foreign's wage. % figdata: dfs_eq
The diagram does comparative statics at a glance. A larger Foreign labour force
raises $B$: Home makes fewer goods at a higher relative wage. A shift of
spending toward Home's goods raises $\vartheta$ at every $z$, with the same effect.

\begin{definition}[Terms of trade]\label{def:b3:ricardian-trade:terms-of-trade}
\index{terms of trade}
A country's \emph{terms of trade} are the price of its exports relative to the
price of its imports. In the Ricardian model the price of a Home export $z$ is
$w a(z)$ and of a Home import $z'$ is $w^* a^*(z')$, so Home's terms of trade
between them are $\omega\, a(z)/a^*(z')$: given technology, they move one for one
with the equilibrium relative wage.
\end{definition}

\section{Transport costs and non-traded goods}\label{sec:b3:ricardian-trade:iceberg}

\begin{definition}[Iceberg trade cost; non-traded good]\label{def:b3:ricardian-trade:iceberg}
\index{iceberg trade cost}\index{non-traded good}
A trade cost is of the \emph{iceberg} form if, for one unit of a good to arrive
abroad, $\kappa \ge 1$ units must be shipped; the rest ``melts'' on the way. A
\emph{non-traded good} is one that each country makes for itself, because the
cost of shipping exceeds the difference in production costs.
\end{definition}

\begin{proposition}[A band of non-traded goods]\label{prop:b3:ricardian-trade:band}
With \omterm{def:b3:ricardian-trade:iceberg}{iceberg} factor $\kappa > 1$ and relative wage $\omega$, Home exports every good
with $A(z) > \kappa\omega$, imports every good with $A(z) < \omega/\kappa$, and both
countries make for themselves every good with $\omega/\kappa < A(z) < \kappa\omega$.
\end{proposition}

\begin{proof}
Delivered abroad, Home's good costs $\kappa w a(z)$, against Foreign's own cost
$w^* a^*(z)$: Home sells abroad if $\kappa w a(z) < w^* a^*(z)$, that is $A(z) >
\kappa\omega$. At Home, Foreign's good costs $\kappa w^* a^*(z)$ against $w a(z)$:
Foreign sells at Home if $A(z) < \omega/\kappa$. Between the two bounds each
country's own good is cheaper in its own market. Since $\kappa > 1$, $\omega/\kappa <
\kappa\omega$ and the band is not empty.
\end{proof}

\begin{omfigure}
\begin{tikzpicture}[econ, xscale=7, yscale=1.1]
  \draw[econ axes] (0,0) -- (1.05,0) node[right] {$z$};
  \draw[econ axes] (0,0) -- (0,3.6) node[above] {$A(z)$};
  \draw[curve, omDef, domain=0.02:1, samples=50] plot (\x, {3.2*exp(-2.3*\x)});
  \draw[guide] (0,1.65) node[left] {$\kappa\omega$} -- (1,1.65);
  \draw[guide] (0,1.0) node[left] {$\omega/\kappa$} -- (1,1.0);
  \draw[guide] (0.288,0) -- (0.288,1.65);
  \draw[guide] (0.506,0) -- (0.506,1.0);
  \node[font=\scriptsize, below] at (0.144,-0.05) {Home exports};
  \node[font=\scriptsize, below, align=center] at (0.397,-0.05) {non-traded};
  \node[font=\scriptsize, below] at (0.76,-0.05) {Home imports};
\end{tikzpicture}
\par\smallskip{\small With an \omterm{def:b3:ricardian-trade:iceberg}{iceberg} cost, the relative wage cuts the chain
twice: goods between $\omega/\kappa$ and $\kappa\omega$ are not traded. A fall in
$\kappa$ narrows the band, and goods at both of its ends start to be traded.}
\end{omfigure}

The band explains a fact the frictionless model cannot: most of what a
country consumes --- haircuts, housing, meals, much of construction --- is made
at home, in every country, rich or poor. Services with high trade costs are the
\omterm{def:b3:ricardian-trade:iceberg}{non-traded goods} of the model.

\section{A productivity gain abroad}\label{sec:b3:ricardian-trade:gain}

When a trading partner becomes more productive, is Home threatened? The model
gives a sharp answer for a gain spread evenly over all goods.

\begin{proposition}[A uniform productivity gain abroad raises Home's real wage]\label{prop:b3:ricardian-trade:abroad}
In the DFS model, let Foreign's labour requirements all fall by the same factor,
$a^*(z) \to a^*(z)/(1 + \gamma)$ with $\gamma > 0$. Then Home's real wage in terms of
every good rises or stays the same.
\end{proposition}

\begin{proof}
The new $A$ is $A(z)/(1 + \gamma)$, lower at every $z$, while $B$ is unchanged; the
curves now cross at $\bar z' \le \bar z$, with $\omega' = A(\bar z')/(1 + \gamma)$. For a
good that Home still makes, its real wage in that good, $w/(w a(z)) = 1/a(z)$, is
unchanged. For a good that Home imports after the gain, its real wage is
$w/(w^* a^*(z)(1 + \gamma)^{-1}) = \omega'(1 + \gamma)/a^*(z) = A(\bar z')/a^*(z)$, and
$A(\bar z') \ge A(\bar z) = \omega$ because $A$ is decreasing and $\bar z' \le \bar z$; so it
is at least $\omega/a^*(z)$, the old real wage in that good if it was imported, and
at least $1/a(z)$ if Home made it before, since it is now cheaper to import.
\end{proof}

Foreign's gain lowers Home's relative wage but lowers the prices of Home's
imports by more. A gain concentrated in the goods Home exports is another matter:
it can erode Home's \omterm{def:b3:ricardian-trade:comparative-advantage}{comparative advantage} and its gains from trade
(\cref{exo:b3:ricardian-trade:6}).

\section{Many countries and the gains from trade}\label{sec:b3:ricardian-trade:many}

With many countries, each good is bought from the cheapest supplier after trade
costs, and the pattern is no longer a single chain. Eaton and Kortum (2002)
made the model tractable by assuming that each country's productivity in each
good is a random draw from one family of distributions. Bilateral trade then
falls with distance in a way that can be estimated, and the response of trade to
trade costs is governed by one number, the \emph{trade elasticity} $\varepsilon$.
Arkolakis, Costinot and Rodríguez-Clare (2012) showed that in this model, and in a
wide class of others, the gains from trade depend on two statistics only.

\begin{definition}[Domestic expenditure share; gains-from-trade formula]\label{def:b3:ricardian-trade:acr}
\index{domestic expenditure share}\index{gains-from-trade formula}
A country's \emph{domestic expenditure share} $\lambda$ is the share of its
spending on goods made at home. The \emph{gains-from-trade formula} gives the
change in real income $W$ after a change in trade costs as $\widehat W =
\hat\lambda^{-1/\varepsilon}$, where a hat is the ratio of new to old values.
\end{definition}

\begin{proposition}[The ACR formula]\label{prop:b3:ricardian-trade:acr}
In models with one factor, perfect or monopolistic competition, balanced trade
and a constant trade elasticity $\varepsilon$, the formula of
\cref{def:b3:ricardian-trade:acr} holds. In particular, moving to autarky
($\lambda' = 1$) multiplies real income by $\lambda^{1/\varepsilon}$, so the gain from
trade over autarky is $\lambda^{-1/\varepsilon} - 1$.
\end{proposition}

\admitted
The proof expresses the price index through the cost of the domestic good and
its share in spending; it is in Arkolakis, Costinot and Rodríguez-Clare (2012).

\begin{measure}[Domestic expenditure share]
$\lambda$ is measured as one minus the share of imports in a country's absorption,
its spending on goods and services: absorption is output minus exports plus
imports, with output measured gross (including inputs) when the imports include
intermediate goods. A large economy that trades little relative to its size has
$\lambda$ near 0.9; a small open one, much lower. The formula's other input, $\varepsilon$,
comes from estimates of how trade responds to trade costs. Both matter: with
$\lambda = 0.8$, an elasticity of 5 gives a gain of $0.8^{-1/5} - 1 = 4.6~\%$, and an
elasticity of 10 half that.
\end{measure}

\section{The tests}\label{sec:b3:ricardian-trade:tests}

The model's central prediction is about \omterm{def:b3:ricardian-trade:autarky}{autarky prices}: a country exports the
goods that were relatively cheap before it traded. A country's \omterm{def:b3:ricardian-trade:autarky}{autarky prices}
are almost never observed, because almost every country has always traded.
Japan before 1859 is the exception, and Bernhofen and Brown used it to test a
general form of the prediction.

\begin{proposition}[The law of comparative advantage]\label{prop:b3:ricardian-trade:law}
Suppose a country's production maximises the value of output at the prices it
faces, its consumers choose the best bundle they can afford, and with trade they
consume a bundle $C$ they could not have chosen in autarky. Let $T = Q - C$ be its
net exports with trade, $Q$ its output with trade and $p^A$ its \omterm{def:b3:ricardian-trade:autarky}{autarky prices}.
Then $p^A \cdot T < 0$: its trade, valued at \omterm{def:b3:ricardian-trade:autarky}{autarky prices}, is negative.
\end{proposition}

\begin{proof}
In autarky consumption equals output $Q^A$, and since $C$ was not chosen, it
cost more at \omterm{def:b3:ricardian-trade:autarky}{autarky prices}: $p^A \cdot C > p^A \cdot Q^A$. Production in autarky
maximised the value of output at $p^A$ over all feasible outputs, so $p^A \cdot Q^A
\ge p^A \cdot Q$. Together, $p^A \cdot C > p^A \cdot Q$, that is $p^A \cdot (Q - C) < 0$.
\end{proof}

The proposition needs no assumption about the number of goods or the
technology; it is the many-good Ricardian prediction stated so that it can be
checked with one vector of prices and one of trade flows. The opening chart is
the check, year by year.

\begin{record}[Japan opens its ports, 1859]
Bernhofen and Brown collected Japanese prices of 1851--53 for the goods Japan
later traded, and trade flows from 1868, when comprehensive statistics begin.
Silk and silkworm eggs made up 51.6~\% of exports in 1868--75 and tea 28.2~\%;
cotton yarn and cloth 33.5~\% of imports and woollens 19.2~\%. % ledger: b3rt:bb04.t1
Valued at \omterm{def:b3:ricardian-trade:autarky}{autarky prices}, net trade was negative in every year from 1868 to 1875.
% ledger: b3rt:bb04.t2.1868, b3rt:bb04.t2.1875
In a second paper they asked how much richer Japan would have had to be in
1851--53 to afford the bundle trade would have offered, and put the upper bound
at about 8 to 9~\% of GDP. % ledger: b3rt:bb05.gains
\end{record}

\begin{record}[Labour productivity and exports]
Tests on modern data compare countries' exports sector by sector with their
relative labour productivity. Golub and Hsieh, for several countries against
the United States, found the coefficients on relative productivity and unit
labour costs almost always of the predicted sign and statistically significant,
while much of the variation of trade across sectors remained unexplained.
% ledger: b3rt:gh00.result
\end{record}

A sharper test needs the productivity of every country in every good,
including goods it does not make. Agriculture provides it: agronomists can
compute what a field would yield in a crop that is not planted. Costinot and
Donaldson used such potential yields for 17 crops on 1.6 million parcels of land
in 55 countries, and found that Ricardo's theory has significant explanatory
power for what is grown. % ledger: b3rt:cd12.data
The chart shows why actual yields could not have done the job.

\begin{omchart}{Relative yields and relative areas, wheat against maize,
2020--2022 averages, 73 countries with at least 10{,}000 ha of each crop in each
year. These are actual yields on the land planted, not the potential yields of
the published test. Line: least squares.}{FAOSTAT (crops); our computation}
\begin{tikzpicture}
\begin{axis}[omscatter, width=0.72\linewidth, height=0.45\linewidth, xmin=-2, xmax=2,
  ymin=-6, ymax=6, xlabel={$\ln$(wheat yield / maize yield)},
  ylabel={$\ln$(wheat area / maize area)},
  x tick label style={/pgf/number format/.cd, fixed, precision=1}]
  \addplot table[x=x, y=y] {figdata/bachelor-3/07-ricardian-trade/out/crops.csv};
  \addplot[omThm, thick, no marks, sharp plot] table[x=x, y=y] {figdata/bachelor-3/07-ricardian-trade/out/crops_fit.csv};
\end{axis}
\end{tikzpicture}
\end{omchart}

The relation is negative: the correlation is $-0.30$. % figdata: crops_stats
Countries that plant mostly wheat, such as Kazakhstan, Iran and Australia, have
high maize yields relative to wheat, because the little maize they grow is on
their best or irrigated land; countries that plant mostly maize do the same with
wheat. The yield of a crop on the land where it is grown is chosen along with
the land, so it measures a selected sample of fields. Potential yields, computed
for every field and every crop, remove the selection, and with them the
prediction holds.

\begin{model}{The Ricardian model with many goods}
\assumes One factor of production, labour, with constant labour requirements that
differ across countries and goods; perfect competition; balanced trade; full
employment.
\predicts Each country exports the goods in which its relative labour cost is
lowest, cut along the \omterm{def:b3:ricardian-trade:comparative-advantage}{chain of comparative advantage} by the relative wage; net
trade valued at \omterm{def:b3:ricardian-trade:autarky}{autarky prices} is negative; trade raises real wages in both
countries; a uniform productivity gain abroad raises Home's real wage.
\tested Japan after 1859: net trade at \omterm{def:b3:ricardian-trade:autarky}{autarky prices} negative in each of the
eight years observed. Sectoral exports rise with relative labour productivity
across countries; potential crop yields predict what farms grow.
\breaks Complete specialisation, which the model predicts for almost every good,
is rare: countries make and export goods in the same industries. Capital, land
and skills differ across countries (\cref{ch:b3:factor-endowments}), and firms
within an industry differ (\cref{ch:b3:scale-variety-firms}).
\end{model}

\section{Exercises}\label{sec:b3:ricardian-trade:exercises}

\begin{exercise}[$\star$]\label{exo:b3:ricardian-trade:1}
In the opening chart, what is the sign of the total in each year, and what does
the law of \omterm{def:b3:ricardian-trade:comparative-advantage}{comparative advantage} predict? In which year is the total closest to
zero, and why does that year still count as a success?
\end{exercise}

\begin{exercise}[$\star$]\label{exo:b3:ricardian-trade:2}
Five goods need 2, 3, 4, 6 and 5 hours at Home and 10, 9, 6, 6 and 3 hours abroad.
Rank them on the \omterm{def:b3:ricardian-trade:comparative-advantage}{chain of comparative advantage}. Which goods does Home make at a
relative wage of 1.2? Of 3.5?
\end{exercise}

\begin{exercise}[$\star$]\label{exo:b3:ricardian-trade:3}
In Brenland a unit of cloth needs 4 hours and of grain 2; in Kestria 6 and 12.
\begin{enumerate}
\item Give each country's \omterm{def:b3:ricardian-trade:autarky}{autarky price} of cloth in grain.
\item After trade, Brenland makes only grain, Kestria only cloth, and the relative
wage is $\omega = 2$. Compute the world price of cloth in grain, and Brenland's
\omterm{def:b3:ricardian-trade:terms-of-trade}{terms of trade}.
\end{enumerate}
\end{exercise}

\begin{exercise}[$\star$]\label{exo:b3:ricardian-trade:4}
What pattern in Japan's prices and trade after 1859 would have contradicted the
law of \omterm{def:b3:ricardian-trade:comparative-advantage}{comparative advantage}?
\end{exercise}

\begin{exercise}[$\star\star$]\label{exo:b3:ricardian-trade:5}
In a DFS model $A(z) = 3 - 2z$, spending shares are equal ($\vartheta(z) = z$) and
$L = L^*$. Find $\bar z$ and $\omega$.
\end{exercise}

\begin{exercise}[$\star\star$]\label{exo:b3:ricardian-trade:6}
Take the two goods of \cref{exo:b3:ricardian-trade:3}, equal spending shares and
a Kestrian labour force twice Brenland's.
\begin{enumerate}
\item Check that $\omega = 2$ with complete specialisation is the equilibrium, and
compute Brenland's gain over autarky in real income.
\item Kestria's grain now needs 6 hours. What changes for Brenland?
\item Kestria's grain now needs 3 hours. What happens to trade and to Brenland's
real income? Compare with \cref{prop:b3:ricardian-trade:abroad}.
\end{enumerate}
\end{exercise}

\begin{exercise}[$\star\star$]\label{exo:b3:ricardian-trade:7}
For the five goods of \cref{exo:b3:ricardian-trade:2}, at a relative wage of 1.3
and an \omterm{def:b3:ricardian-trade:iceberg}{iceberg} factor of 1.2, which goods does Home export, import and make for
itself?
\end{exercise}

\begin{exercise}[$\star\star$]\label{exo:b3:ricardian-trade:8}
Using FAOSTAT through the series loader, compute the correlation between
relative yields and relative areas for wheat against maize, as in the chart,
and for wheat against barley. Why might the two differ?
\end{exercise}

\begin{exercise}[$\star\star$]\label{exo:b3:ricardian-trade:9}
A country's \omterm{def:b3:ricardian-trade:acr}{domestic expenditure share} is 0.8. Compute its gain from trade over
autarky with a trade elasticity of 5 and of 10. Which input would you most like
to know better?
\end{exercise}

\begin{exercise}[$\star\star\star$]\label{exo:b3:ricardian-trade:10}
In the DFS model with an \omterm{def:b3:ricardian-trade:iceberg}{iceberg} factor $\kappa > 1$, let $z_E$ and $z_M$ solve
$A(z_E) = \kappa\omega$ and $A(z_M) = \omega/\kappa$. Show that Home's wage bill equals its
earnings when $\omega L(1 - \vartheta(z_M)) = \vartheta(z_E) L^*$, and that this equation has
exactly one solution $\omega$.
\end{exercise}

\begin{exercise}[$\star\star\star$]\label{exo:b3:ricardian-trade:11}
Trade costs fall, from $\kappa = 1.2$ to $\kappa = 1.1$, with $\omega$ held at its old
value. Which goods start to be traded? Then let $\omega$ adjust: using
\cref{exo:b3:ricardian-trade:10}, in which direction does it move if the goods at
the two ends of the band carry equal spending shares and $L = L^*$?
\end{exercise}

\begin{exercise}[$\star\star\star$]\label{exo:b3:ricardian-trade:12}
Countries export and import goods in the same industries, and few industries
are left entirely to foreign producers. Give two reasons inside the Ricardian
model and two outside it.
\end{exercise}

\section{Problem: Brenland and Kestria with Twelve Goods}\label{sec:b3:ricardian-trade:problem}

\begin{problem}[{Weekend problem --- Brenland and Kestria with twelve goods}]\label{pb:b3:ricardian-trade:twelve}
Brenland and Kestria (invented) make twelve goods with labour alone. The table
gives the hours of labour per unit in each country; cloth and grain are the
series' usual values. Consumers in both countries spend one twelfth of their
income on each good, and the two labour forces are equal. % figdata: twelve_goods

\begin{center}\small
\pgfplotstabletypeset[col sep=comma, skip rows between index={6}{12},
  columns/good/.style={column name={Good}, string type, column type=l},
  columns/hours_B/.style={column name={Brenland}, fixed, precision=1},
  columns/hours_K/.style={column name={Kestria}, fixed, precision=1},
  every head row/.style={after row=\hline}, every last row/.style={after row=\hline}]
  {figdata/bachelor-3/07-ricardian-trade/out/twelve_goods.csv}
\qquad
\pgfplotstabletypeset[col sep=comma, skip rows between index={0}{6},
  columns/good/.style={column name={Good}, string type, column type=l},
  columns/hours_B/.style={column name={Brenland}, fixed, precision=1},
  columns/hours_K/.style={column name={Kestria}, fixed, precision=1},
  every head row/.style={after row=\hline}, every last row/.style={after row=\hline}]
  {figdata/bachelor-3/07-ricardian-trade/out/twelve_goods.csv}
\end{center}

\medskip\noindent\textbf{Part I --- autarky and the chain.}
\begin{enumerate}
\item Compute $A_i$ for each good and check that the table lists the goods in
the order of the chain.
\item In which goods does Brenland have an absolute advantage? Does it matter
for the pattern of trade?
\item Give each country's \omterm{def:b3:ricardian-trade:autarky}{autarky price} of cloth in grain, and of radios in
grain.
\item Which goods would Brenland make at $\omega = 1$? At $\omega = 2$?
\item Explain why the equilibrium $\omega$ must lie between 0.6 and 6.
\item Show that with equal spending shares and equal labour forces, trade
balances when $\omega$ equals the number of goods Brenland makes divided by the
number Kestria makes.
\end{enumerate}

\medskip\noindent\textbf{Part II --- the free-trade equilibrium.}
\begin{enumerate}[resume]
\item Show that $\omega = 1.4$, with Brenland making the first seven goods, is an
equilibrium.
\item Why is $\omega = 1.2$ not an equilibrium? And $\omega = 2$?
\item Compute Brenland's \omterm{def:b3:ricardian-trade:terms-of-trade}{terms of trade} between grain and radios, in autarky
and with trade.
\item Compute Brenland's real wage in radios, in autarky and with trade.
\item With equal shares, Brenland's real income over autarky is multiplied by
the twelfth root of the product of $\omega/A_i$ over the goods it imports. Explain
why, and compute the gain.
\item Compute Kestria's gain the same way. Which country gains more, and why?
\end{enumerate}

\medskip\noindent\textbf{Part III --- a 10~\% iceberg cost.}
\begin{enumerate}[resume]
\item With $\kappa = 1.1$, state the conditions for Brenland to export, to import
and to make a good for itself.
\item At $\omega = 1.5$, list the goods in each group.
\item Check that $\omega = 1.5$ is the equilibrium.
\item Why does Brenland's relative wage rise from 1.4 to 1.5?
\item Compute Brenland's gain over autarky with the \omterm{def:b3:ricardian-trade:iceberg}{iceberg} cost; Kestria's is
32.8~\%. Compare with Part II.
\end{enumerate}

\medskip\noindent\textbf{Part IV --- a productivity gain in Kestria, and the
formula.}
\begin{enumerate}[resume]
\item Trade is free again, and Kestria's labour requirements all fall by a sixth
(its productivity rises by 20~\%). Give the new values of $A_i$.
\item Show that the new equilibrium has $\omega = 1.25$, with cloth made in both
countries, and find Brenland's share of the world's cloth.
\item Brenland's gain over autarky is now 20.3~\%. Check the direction against
\cref{prop:b3:ricardian-trade:abroad}.
\item Kestria's gain over its own new autarky falls from 45.0~\% to 39.2~\%. Has
Kestria become worse off?
\item Compute Brenland's \omterm{def:b3:ricardian-trade:acr}{domestic expenditure share} in free trade (Part II) and
the gain given by the formula of \cref{prop:b3:ricardian-trade:acr} with $\varepsilon =
5$.
\item What trade elasticity would make the formula give the gain computed in
Part II?
\item Why do the two methods differ in this model?
\item State the results: Brenland's gain from trade, in percent of real income,
computed two ways.
\end{enumerate}
\end{problem}
